\documentclass[aspectratio=169]{beamer}

% Tema UFS --- Universidade Federal de Sergipe
\usetheme{UFS}

\usepackage[utf8]{inputenc}
\usepackage[portuguese]{babel}
\usepackage{amsmath, amssymb, amsthm}
\usepackage{graphicx}
\usepackage{booktabs}
\usepackage{xcolor}

\usepackage{tikz}
\usetikzlibrary{shapes.geometric, arrows, arrows.meta, positioning, matrix, fit,
  backgrounds, decorations.pathreplacing, shadows, patterns,
  automata, intersections, calc, shapes}

% ---------- Listings: paleta VS Code Light+ (estilo UFS) ----------
\usepackage{listings}
\definecolor{bgcolor}{HTML}{FFFFFF}
\definecolor{linecolor}{HTML}{DDDDDD}
\definecolor{numcolor}{HTML}{999999}
\definecolor{kwcolor}{HTML}{0000FF}
\definecolor{strcolor}{HTML}{A31515}
\definecolor{cmtcolor}{HTML}{008000}

\lstdefinestyle{ufsStyle}{
  backgroundcolor=\color{bgcolor},
  basicstyle=\ttfamily\scriptsize\color{black},
  keywordstyle=\color{kwcolor}\bfseries,
  stringstyle=\color{strcolor},
  commentstyle=\color{cmtcolor}\itshape,
  numberstyle=\tiny\color{numcolor},
  numbers=left, numbersep=8pt,
  xleftmargin=14pt, framexleftmargin=14pt,
  breaklines=true, tabsize=2,
  keepspaces=true, showspaces=false,
  showstringspaces=false, showtabs=false,
  frame=single, rulecolor=\color{linecolor},
  framesep=3pt, captionpos=b, upquote=true,
  columns=fullflexible,
  morekeywords={var, record, sealed, permits},
}
\lstset{style=ufsStyle, language=Java}

% ---------- TikZ styles para árvores (cores UFS) ----------
\tikzset{
  bstnode/.style={
    circle, draw=UFSBlue, fill=UFSBlue!15, thick,
    minimum size=7mm, inner sep=0pt, font=\small\bfseries
  },
  hl/.style={fill=UFSYellow!60, draw=UFSYellow!80!black},
  nilnode/.style={
    rectangle, draw=UFSGray, fill=UFSLightGray,
    minimum size=4mm, inner sep=1pt, font=\tiny
  },
  edge/.style={-, thick, UFSBlue}
}

\title[BST]{Árvores Binárias de Busca (BST)}
\subtitle{Aula 08 --- COMP0497}
\author{Prof.\ André Yoshiaki Kashiwabara}
\institute{DCOMP --- Universidade Federal de Sergipe\\
Algoritmos e Estruturas de Dados 1 (COMP0497)}
\date{\today}

\begin{document}

% ============================================================
\begin{frame}[plain]
  \titlepage
\end{frame}

% ============================================================
% ============================================================
\begin{frame}{O que você vai aprender hoje}
\begin{center}
\begin{tikzpicture}[
  item/.style={rectangle, draw=UFSBlue, fill=UFSBlue!10, rounded corners,
    minimum width=7.5cm, minimum height=0.55cm, font=\small}
]
  \node[item] (t1) at (0, 0.0)  {1. Árvores binárias: conceitos e representação};
  \node[item] (t2) at (0,-0.8)  {2. Travessias em profundidade (in/pre/post-order)};
  \node[item] (t3) at (0,-1.6)  {3. Travessia em largura (BFS)};
  \node[item] (t4) at (0,-2.4)  {4. BST: propriedade, busca, mínimo, sucessor};
  \node[item] (t5) at (0,-3.2)  {5. BST: inserção e remoção};
  \node[item] (t6) at (0,-4.0)  {6. Análise: altura $h$ e por que balancear};
  \node[item] (t7) at (0,-4.8)  {7. Exercícios em sala};
  \draw[->, thick] (t1)--(t2); \draw[->, thick] (t2)--(t3);
  \draw[->, thick] (t3)--(t4); \draw[->, thick] (t4)--(t5);
  \draw[->, thick] (t5)--(t6); \draw[->, thick] (t6)--(t7);
\end{tikzpicture}
\end{center}
\end{frame}

% ============================================================
\section{Árvores binárias: conceitos e representação}

\begin{frame}[shrink=10]{O que é uma árvore binária?}
\begin{block}{Definição}
Uma \textbf{árvore binária} é um conjunto finito de nós que é
\emph{vazio}, ou consiste em um nó \emph{raiz} mais duas árvores
binárias disjuntas chamadas \emph{subárvore esquerda} e
\emph{subárvore direita}.
\end{block}
\begin{columns}[T]
\begin{column}{0.55\textwidth}
\textbf{Terminologia:}
\begin{itemize}\small
  \item \textbf{Raiz:} único nó sem pai.
  \item \textbf{Filho / pai:} relação entre nós ligados.
  \item \textbf{Folha:} nó sem filhos.
  \item \textbf{Nó interno:} tem pelo menos um filho.
  \item \textbf{Irmãos:} filhos do mesmo pai.
  \item \textbf{Profundidade} de $x$: nº de arestas até a raiz.
  \item \textbf{Nível $k$:} todos os nós de profundidade $k$.
  \item \textbf{Altura} $h$: maior profundidade da árvore.
\end{itemize}
\end{column}
\begin{column}{0.45\textwidth}
\centering
\begin{tikzpicture}[level distance=8mm, sibling distance=18mm,
  level 2/.style={sibling distance=10mm}]
  \node[bstnode,label=right:{\scriptsize raiz, prof. 0}]{A}
    child{ node[bstnode]{B}
      child{ node[bstnode]{D}} child{ node[bstnode]{E}} }
    child{ node[bstnode]{C}
      child[missing]
      child{ node[bstnode,label=right:{\scriptsize folha, prof. 2}]{F}} };
\end{tikzpicture}\\[2pt]
\small $h = 2$, 6 nós, 3 folhas (D, E, F).
\end{column}
\end{columns}
\end{frame}

\begin{frame}[shrink=15]{Tipos especiais de árvores binárias}
\begin{columns}[T]
\begin{column}{0.33\textwidth}
\centering\textbf{Cheia (full)}\\[2pt]
\begin{tikzpicture}[level distance=7mm, sibling distance=14mm,
  level 2/.style={sibling distance=8mm}]
  \node[bstnode]{}
    child{ node[bstnode]{}
      child{ node[bstnode]{}} child{ node[bstnode]{}} }
    child{ node[bstnode]{}};
\end{tikzpicture}\\[2pt]
\scriptsize Todo nó tem 0 ou 2 filhos.
\end{column}
\begin{column}{0.33\textwidth}
\centering\textbf{Completa}\\[2pt]
\begin{tikzpicture}[level distance=7mm, sibling distance=14mm,
  level 2/.style={sibling distance=8mm}]
  \node[bstnode]{}
    child{ node[bstnode]{}
      child{ node[bstnode]{}} child{ node[bstnode]{}} }
    child{ node[bstnode]{}
      child{ node[bstnode]{}} child[missing] };
\end{tikzpicture}\\[2pt]
\scriptsize Todos os níveis cheios, exceto talvez o último (preenchido à esquerda).
\end{column}
\begin{column}{0.33\textwidth}
\centering\textbf{Perfeita}\\[2pt]
\begin{tikzpicture}[level distance=7mm, sibling distance=14mm,
  level 2/.style={sibling distance=8mm}]
  \node[bstnode]{}
    child{ node[bstnode]{}
      child{ node[bstnode]{}} child{ node[bstnode]{}} }
    child{ node[bstnode]{}
      child{ node[bstnode]{}} child{ node[bstnode]{}} };
\end{tikzpicture}\\[2pt]
\scriptsize Todos os níveis completamente cheios.
\end{column}
\end{columns}
\vspace{6pt}
\begin{block}{Limites de altura}
Para uma árvore binária com $n$ nós:
\[
\lceil \log_2(n+1) \rceil - 1 \;\le\; h \;\le\; n - 1
\]
O \emph{mínimo} é atingido por árvores perfeitas; o \emph{máximo}, por cadeias (uma sequência linear).
\end{block}
\end{frame}

\begin{frame}[fragile]{Representação ligada}
\begin{columns}[T]
\begin{column}{0.45\textwidth}
\centering
\begin{tikzpicture}[
  node distance=12mm,
  cell/.style={rectangle split, rectangle split parts=3,
    rectangle split horizontal, draw=UFSBlue, fill=UFSBlue!10,
    minimum height=6mm, font=\small\ttfamily},
  null/.style={rectangle, draw=UFSGray, fill=UFSLightGray,
    minimum size=3mm, inner sep=1pt, font=\tiny}
]
  \node[cell] (r) {$\bullet$\nodepart{two}5\nodepart{three}$\bullet$};
  \node[cell, below left=10mm and 0mm of r] (a) {$\bullet$\nodepart{two}3\nodepart{three}$\bullet$};
  \node[cell, below right=10mm and 0mm of r] (b) {nil\nodepart{two}8\nodepart{three}nil};
  \node[null, below left =6mm and -2mm of a] (n1) {nil};
  \node[null, below right=6mm and -2mm of a] (n2) {nil};
  \draw[->, UFSBlue] (r.one south) -- (a.north);
  \draw[->, UFSBlue] (r.three south) -- (b.north);
  \draw[->, UFSBlue] (a.one south) -- (n1.north);
  \draw[->, UFSBlue] (a.three south) -- (n2.north);
\end{tikzpicture}\\[2pt]
\scriptsize Cada célula: \texttt{left | key | right}
\end{column}
\begin{column}{0.55\textwidth}
\begin{lstlisting}
class Node<K, V> {
    K key;
    V value;
    Node<K, V> left, right;   // filhos
    Node<K, V> parent;        // opcional (CLRS)

    Node(K key, V value) {
        this.key = key;
        this.value = value;
    }
}
\end{lstlisting}
\begin{itemize}\scriptsize
  \item \textbf{Sem} \texttt{parent}: estilo Sedgewick. Suficiente para BST.
  \item \textbf{Com} \texttt{parent}: estilo CLRS. Facilita sucessor e remoção iterativos.
\end{itemize}
\end{column}
\end{columns}
\end{frame}

\begin{frame}{Representação por array (alternativa)}
\begin{columns}[T]
\begin{column}{0.55\textwidth}
Para árvores \emph{completas} (heaps), pode-se guardar a árvore em um \textbf{vetor} sem ponteiros:
\begin{itemize}
  \item raiz no índice 0;
  \item filhos de $i$: $2i+1$ (esq) e $2i+2$ (dir);
  \item pai de $i$: $\lfloor (i-1)/2 \rfloor$.
\end{itemize}
\begin{exampleblock}{Quando usar?}
Ótimo para \alert{heaps} e árvores balanceadas estáticas.
\emph{Ruim} para BSTs gerais — desperdiça memória se a árvore for esparsa.
\end{exampleblock}
\end{column}
\begin{column}{0.45\textwidth}
\centering
\begin{tikzpicture}[level distance=8mm, sibling distance=18mm,
  level 2/.style={sibling distance=9mm}]
  \node[bstnode]{A}
    child{ node[bstnode]{B}
      child{ node[bstnode]{D}} child{ node[bstnode]{E}} }
    child{ node[bstnode]{C}
      child{ node[bstnode]{F}} child{ node[bstnode]{G}} };
\end{tikzpicture}\\[6pt]
{\small\ttfamily
\begin{tabular}{|c|c|c|c|c|c|c|}
\hline
A & B & C & D & E & F & G \\
\hline
0 & 1 & 2 & 3 & 4 & 5 & 6 \\
\hline
\end{tabular}}
\end{column}
\end{columns}
\end{frame}

% ============================================================
\section{Definição e propriedade da BST}

\begin{frame}{Por que precisamos de BST?}
\begin{columns}[T]
\begin{column}{0.5\textwidth}
\textbf{Problema:}\\[2pt]
Manter um \emph{dicionário ordenado} com inserções,
remoções e buscas frequentes.

\begin{itemize}
  \item \alert{Vetor ordenado:} busca $O(\log n)$ \alert{mas} inserção $O(n)$.
  \item \alert{Lista encadeada:} inserção $O(1)$ \alert{mas} busca $O(n)$.
\end{itemize}
\end{column}
\begin{column}{0.5\textwidth}
\textbf{Solução: BST}\\[2pt]
Estrutura ligada que mantém a ordem
e oferece todas as operações em \emph{tempo proporcional à altura}~$h$.
\begin{itemize}
  \item Inserção, busca, remoção: $O(h)$.
  \item Em árvore balanceada: $h = O(\log n)$.
\end{itemize}
\end{column}
\end{columns}
\vspace{6pt}
\begin{block}{Suporta}
\textsc{Search, Min, Max, Predecessor, Successor, Insert, Delete}
\hfill\tiny{(Cormen et al., cap. 12)}
\end{block}
\end{frame}

\begin{frame}{BST: Definição}
\begin{block}{Propriedade da árvore binária de busca}
Para todo nó $x$:
\begin{itemize}
  \item se $y$ está na \emph{subárvore esquerda} de $x$: $\mathit{key}[y] \le \mathit{key}[x]$;
  \item se $y$ está na \emph{subárvore direita}  de $x$: $\mathit{key}[y] \ge \mathit{key}[x]$.
\end{itemize}
\end{block}
\begin{exampleblock}{Intuição}
``Tudo à esquerda é menor (ou igual), tudo à direita é maior (ou igual).''
Esta invariante vale \emph{recursivamente} em cada subárvore.
\end{exampleblock}
\end{frame}

\begin{frame}{BST: exemplo visual}
\begin{columns}[T]
\begin{column}{0.5\textwidth}
\centering\textbf{(a) altura 2 --- eficiente}\\[4pt]
\begin{tikzpicture}[level distance=10mm, sibling distance=18mm,
  level 2/.style={sibling distance=9mm}]
  \node[bstnode]{5}
    child{ node[bstnode]{3}
      child{ node[bstnode]{2}} child{ node[bstnode]{5}} }
    child{ node[bstnode]{7}
      child[missing] child{ node[bstnode]{8}} };
\end{tikzpicture}
\end{column}
\begin{column}{0.5\textwidth}
\centering\textbf{(b) altura 4 --- ineficiente}\\[4pt]
\begin{tikzpicture}[level distance=8mm, sibling distance=12mm]
  \node[bstnode]{2}
    child[missing]
    child{ node[bstnode]{3}
      child[missing]
      child{ node[bstnode]{7}
        child{ node[bstnode]{5}
          child[missing] child{ node[bstnode]{5}} }
        child{ node[bstnode]{8}} } };
\end{tikzpicture}
\end{column}
\end{columns}
\vspace{4pt}
\begin{center}\small As mesmas chaves \{2,3,5,5,7,8\}, BSTs diferentes.\\
O custo das operações depende de $h$, não de $n$.\end{center}
\hfill\tiny{Fig.\ 12.1, Cormen et al.}
\end{frame}

% ============================================================
\section{Travessias}

\begin{frame}{Por que travessias?}
\begin{columns}[T]
\begin{column}{0.5\textwidth}
\textbf{Problema:}\\
Como visitar todos os nós em uma ordem útil?
\end{column}
\begin{column}{0.5\textwidth}
\textbf{Solução:} três travessias recursivas naturais.
\begin{itemize}
  \item \textbf{In-order:} esq, \emph{raiz}, dir \;$\Rightarrow$ ordem crescente.
  \item \textbf{Pre-order:} raiz, esq, dir.
  \item \textbf{Post-order:} esq, dir, raiz.
\end{itemize}
\end{column}
\end{columns}
\end{frame}

\begin{frame}[fragile]{In-order: imprime em ordem}
\begin{columns}[T]
\begin{column}{0.45\textwidth}
\centering
\begin{tikzpicture}[level distance=9mm, sibling distance=18mm,
  level 2/.style={sibling distance=10mm}]
  \node[bstnode]{5}
    child{ node[bstnode]{3}
      child{ node[bstnode]{2}} child{ node[bstnode]{5}} }
    child{ node[bstnode]{7}
      child[missing] child{ node[bstnode]{8}} };
\end{tikzpicture}\\[4pt]
\small Saída: \texttt{2 3 5 5 7 8}
\end{column}
\begin{column}{0.55\textwidth}
\begin{lstlisting}
void inorder(Node<K,V> x) {
    if (x == null) return;
    inorder(x.left);
    System.out.println(x.key);
    inorder(x.right);
}
\end{lstlisting}
\begin{block}{Teorema 12.1 (Cormen)}
\textsc{Inorder-Tree-Walk} executa em $\Theta(n)$ em uma BST com $n$ nós.
\end{block}
\end{column}
\end{columns}
\end{frame}

\begin{frame}[fragile,shrink=5]{Pre-order e post-order}
\begin{columns}[T]
\begin{column}{0.38\textwidth}
\centering
\begin{tikzpicture}[level distance=8mm, sibling distance=18mm,
  level 2/.style={sibling distance=10mm}]
  \node[bstnode]{A}
    child{ node[bstnode]{B}
      child{ node[bstnode]{D}} child{ node[bstnode]{E}} }
    child{ node[bstnode]{C}
      child[missing] child{ node[bstnode]{F}} };
\end{tikzpicture}\\[4pt]
\small
\textbf{pre-order:} A B D E C F\\
\textbf{post-order:} D E B F C A
\end{column}
\begin{column}{0.62\textwidth}
\begin{lstlisting}
void preorder(Node<K,V> x) {
    if (x == null) return;
    visit(x);                    // raiz primeiro
    preorder(x.left);
    preorder(x.right);
}

void postorder(Node<K,V> x) {
    if (x == null) return;
    postorder(x.left);
    postorder(x.right);
    visit(x);                    // raiz por ultimo
}
\end{lstlisting}
\begin{itemize}\scriptsize
  \item \textbf{Pre-order:} cópia de árvores, serialização.
  \item \textbf{Post-order:} liberar memória, avaliar expressão (folhas antes da raiz).
\end{itemize}
\end{column}
\end{columns}
\end{frame}

% ============================================================
\section{Travessia em largura (BFS)}

\begin{frame}[shrink=18]{BFS: visitar nível por nível}
\begin{columns}[T]
\begin{column}{0.5\textwidth}
\textbf{Ideia:} usar uma \emph{fila} (FIFO). Visite a raiz, enfileire os filhos, repita.
\begin{itemize}
  \item Visita os nós por \textbf{nível crescente} de profundidade.
  \item Custo: $\Theta(n)$ tempo, $O(\text{largura máxima})$ espaço.
  \item Em árvore perfeita: largura $\approx n/2$.
\end{itemize}
\begin{block}{Quando usar BFS?}
\begin{itemize}\small
  \item Encontrar nó mais próximo da raiz com certa propriedade.
  \item Calcular nível de cada nó.
  \item Imprimir árvore "nível a nível".
\end{itemize}
\end{block}
\end{column}
\begin{column}{0.5\textwidth}
\centering
\begin{tikzpicture}[level distance=8mm, sibling distance=20mm,
  level 2/.style={sibling distance=10mm}]
  \node[bstnode]{1}
    child{ node[bstnode]{2}
      child{ node[bstnode]{4}} child{ node[bstnode]{5}} }
    child{ node[bstnode]{3}
      child{ node[bstnode]{6}} child{ node[bstnode]{7}} };
\end{tikzpicture}\\[4pt]
\small Ordem BFS: \texttt{1 2 3 4 5 6 7}\\
(em contraste: pre-order = 1 2 4 5 3 6 7)
\end{column}
\end{columns}
\end{frame}

\begin{frame}[fragile,shrink=20]{BFS em Java moderno}
\begin{lstlisting}
import java.util.ArrayDeque;
import java.util.Deque;

public void bfs(Node<K, V> root) {
    if (root == null) return;
    Deque<Node<K, V>> fila = new ArrayDeque<>();
    fila.offer(root);                          // enfileira raiz
    while (!fila.isEmpty()) {
        Node<K, V> x = fila.poll();            // remove da frente
        System.out.println(x.key);
        if (x.left  != null) fila.offer(x.left);
        if (x.right != null) fila.offer(x.right);
    }
}
\end{lstlisting}
\begin{exampleblock}{DFS vs BFS}
\textbf{DFS} (in/pre/post-order) usa \emph{pilha} (recursão); explora um ramo de cada vez.\\
\textbf{BFS} usa \emph{fila}; explora a árvore em camadas.
\end{exampleblock}
\end{frame}

% ============================================================
\section{Busca, mínimo, máximo, sucessor}

\begin{frame}[shrink=10]{Busca: o caminho da raiz}
\begin{block}{Ideia}
Compare $k$ com a chave de cada nó. Se $k < \mathit{key}[x]$, vá para a esquerda;
se $k > \mathit{key}[x]$, vá para a direita.
\end{block}
\begin{columns}[T]
\begin{column}{0.45\textwidth}
\centering
\begin{tikzpicture}[level distance=8mm, sibling distance=22mm,
  level 2/.style={sibling distance=11mm},
  level 3/.style={sibling distance=8mm}]
  \node[bstnode,hl]{15}
    child{ node[bstnode,hl]{6}
      child{ node[bstnode]{3} }
      child{ node[bstnode,hl]{7}
        child[missing]
        child{ node[bstnode,hl]{13}
          child{ node[bstnode]{9}} child[missing]
        } } }
    child{ node[bstnode]{18} };
\end{tikzpicture}\\[2pt]
\small Buscando 13: $15 \to 6 \to 7 \to 13$
\end{column}
\begin{column}{0.55\textwidth}
\begin{block}{Custo}
$O(h)$ comparações. Em BST balanceada: $O(\log n)$.
\end{block}
\begin{exampleblock}{Versão iterativa (preferida em Java)}
Evita estouro de pilha em árvores muito altas
e é geralmente mais rápida.
\end{exampleblock}
\end{column}
\end{columns}
\end{frame}

\begin{frame}[fragile]{Busca em Java moderno (genéricos)}
\begin{lstlisting}
public class BST<Key extends Comparable<Key>, Value> {

    public Value get(Key key) {
        Node<Key,Value> x = root;
        while (x != null) {
            int cmp = key.compareTo(x.key);
            if      (cmp < 0) x = x.left;
            else if (cmp > 0) x = x.right;
            else              return x.value;   // achou
        }
        return null;                            // ausente
    }
}
\end{lstlisting}
\begin{itemize}
  \item \texttt{Key extends Comparable<Key>}: aceita qualquer tipo ordenável.
  \item Versão iterativa, $O(h)$ tempo, $O(1)$ espaço extra.
\end{itemize}
\end{frame}

\begin{frame}{Mínimo, máximo, sucessor}
\begin{columns}[T]
\begin{column}{0.5\textwidth}
\textbf{Mínimo:} desça pela esquerda até \texttt{null}.\\[2pt]
\textbf{Máximo:} desça pela direita até \texttt{null}.\\[6pt]
\textbf{Sucessor de $x$ (próxima chave em ordem):}
\begin{itemize}
  \item se $x$ tem subárvore direita: \emph{mínimo} dessa subárvore;
  \item caso contrário: \emph{primeiro ancestral} cujo filho esquerdo também é ancestral de $x$.
\end{itemize}
Tempo: $O(h)$ em ambos os casos.
\end{column}
\begin{column}{0.5\textwidth}
\centering
\begin{tikzpicture}[level distance=8mm, sibling distance=22mm,
  level 2/.style={sibling distance=11mm},
  level 3/.style={sibling distance=8mm}]
  \node[bstnode]{15}
    child{ node[bstnode]{6}
      child{ node[bstnode]{3}
        child{ node[bstnode]{2}} child[missing] }
      child{ node[bstnode]{7}
        child[missing]
        child{ node[bstnode,hl]{13}} } }
    child{ node[bstnode]{18}
      child{ node[bstnode]{17}} child{ node[bstnode]{20}} };
\end{tikzpicture}\\[2pt]
\small Sucessor de 13 é 15 (sobe pelos ancestrais).
\end{column}
\end{columns}
\end{frame}

% ============================================================
\section{Operações ordenadas: rank e select}

\begin{frame}[fragile,shrink=10]{Rank e select com o campo \texttt{size}}
\begin{block}{Por que manter \texttt{size} em cada nó?}
\texttt{x.size = 1 + size(x.left) + size(x.right)} permite consultas \emph{ordinais} em $O(h)$:
\begin{itemize}\small
  \item \textbf{rank(key)}: quantas chaves são \emph{menores} que \texttt{key}?
  \item \textbf{select(k)}: qual é a $k$-ésima menor chave? (índice 0)
\end{itemize}
\end{block}
\begin{columns}[T]
\begin{column}{0.5\textwidth}
\begin{lstlisting}
public int rank(Key key) {
    return rank(root, key);
}
private int rank(Node<Key,Value> x, Key key) {
    if (x == null) return 0;
    int cmp = key.compareTo(x.key);
    if      (cmp < 0) return rank(x.left, key);
    else if (cmp > 0) return 1 + size(x.left)
                          + rank(x.right, key);
    else              return size(x.left);
}
\end{lstlisting}
\end{column}
\begin{column}{0.5\textwidth}
\begin{lstlisting}
public Key select(int k) {
    return select(root, k).key;
}
private Node<Key,Value> select(Node<Key,Value> x, int k) {
    if (x == null) return null;
    int t = size(x.left);
    if      (t > k) return select(x.left,  k);
    else if (t < k) return select(x.right, k - t - 1);
    else            return x;
}
\end{lstlisting}
\end{column}
\end{columns}
\hfill\tiny{Sedgewick \& Wayne, \emph{Algorithms} §3.2 --- também útil para \emph{range queries} em $O(\log n + k)$.}
\end{frame}

% ============================================================
\section{Inserção e remoção}

\begin{frame}[fragile,shrink=15]{Inserção: estilo Sedgewick (recursivo)}
\begin{block}{Ideia}
Buscar pela chave; ao chegar em \texttt{null}, criar nó. A versão recursiva
\emph{retorna a subárvore} e atualiza ponteiros ao subir.
\end{block}
\begin{lstlisting}
public void put(Key key, Value value) {
    root = put(root, key, value);
}
private Node<Key,Value> put(Node<Key,Value> x, Key key, Value value) {
    if (x == null) return new Node<>(key, value);
    int cmp = key.compareTo(x.key);
    if      (cmp < 0) x.left  = put(x.left,  key, value);
    else if (cmp > 0) x.right = put(x.right, key, value);
    else              x.value = value;             // chave ja existe
    x.size = 1 + size(x.left) + size(x.right);     // mantem tamanho
    return x;
}
\end{lstlisting}
Custo: $O(h)$.\hfill\tiny{(Sedgewick, cap.\ 12)}
\end{frame}

\begin{frame}{Inserção: exemplo passo a passo}
\centering
Inserindo \alert{13} em uma BST existente.\\[4pt]
\begin{tikzpicture}[level distance=9mm, sibling distance=20mm,
  level 2/.style={sibling distance=12mm}]
  \node[bstnode]{12}
    child{ node[bstnode]{5}
      child{ node[bstnode]{2}} child{ node[bstnode]{9}} }
    child{ node[bstnode,hl]{18}
      child{ node[bstnode,hl]{15}
        child[edge from parent/.style={dashed,thick,UFSRed}]
          { node[bstnode,hl]{13} } child[missing] }
      child{ node[bstnode]{19}} };
\end{tikzpicture}\\[4pt]
\small O caminho 12 $\to$ 18 $\to$ 15 leva à posição correta. Aresta tracejada é a nova ligação.\\
\hfill\tiny{Fig.\ 12.3, Cormen et al.}
\end{frame}

\begin{frame}[shrink=15]{Remoção: três casos}
Seja $z$ o nó a remover.
\begin{enumerate}
  \item \textbf{$z$ sem filhos:} basta desligar $z$ do pai.
  \item \textbf{$z$ com 1 filho:} pai de $z$ aponta para o único filho de $z$.
  \item \textbf{$z$ com 2 filhos:} substitua $z$ pelo seu \emph{sucessor in-order} $y$
        (o mínimo da subárvore direita); $y$ \emph{tem no máximo 1 filho}.
\end{enumerate}
\vspace{4pt}
\begin{center}
\begin{tikzpicture}[level distance=8mm, sibling distance=14mm,
  level 2/.style={sibling distance=8mm}]
  \node[bstnode]{15}
    child{ node[bstnode]{5}
      child{ node[bstnode]{3}}
      child{ node[bstnode,hl]{12}
        child{ node[bstnode]{10}} child{ node[bstnode]{13}} } }
    child{ node[bstnode]{20}};
\end{tikzpicture}\\
\small Remover \alert{12}: sucessor é 13; copiar 13 para o lugar de 12 e remover o nó 13 original.
\end{center}
\hfill\tiny{Fig.\ 12.4, Cormen et al.\ --- método de Hibbard}
\end{frame}

\begin{frame}[fragile]{Remoção em Java (Hibbard)}
\begin{lstlisting}
private Node<Key,Value> delete(Node<Key,Value> x, Key key) {
    if (x == null) return null;
    int cmp = key.compareTo(x.key);
    if      (cmp < 0) x.left  = delete(x.left,  key);
    else if (cmp > 0) x.right = delete(x.right, key);
    else {
        if (x.right == null) return x.left;     // caso 1 ou 2
        if (x.left  == null) return x.right;    // caso 2
        Node<Key,Value> t = x;                  // caso 3
        x = min(t.right);                       // sucessor
        x.right = deleteMin(t.right);
        x.left  = t.left;
    }
    x.size = 1 + size(x.left) + size(x.right);
    return x;
}
\end{lstlisting}
\end{frame}

% ============================================================
\section{Análise: altura e por que balancear}

\begin{frame}{Altura: o que determina o custo}
\begin{block}{Teorema (Cormen 12.4 / Aslam)}
A altura esperada de uma BST construída \emph{aleatoriamente} por $n$ inserções
de chaves distintas é $\mathbb{E}[h] = O(\log n)$.
\end{block}
\begin{exampleblock}{Mas no pior caso}
Inserindo chaves em ordem crescente, a BST vira uma lista encadeada:
$h = n - 1$ e todas as operações custam $\Theta(n)$.
\end{exampleblock}
\begin{center}
\small
\begin{tabular}{lcc}
\toprule
Operação & BST balanceada & BST degenerada\\
\midrule
\textsc{Search / Insert / Delete} & $O(\log n)$ & $\Theta(n)$\\
\textsc{Min / Max / Successor}    & $O(\log n)$ & $\Theta(n)$\\
\bottomrule
\end{tabular}
\end{center}
\end{frame}

\begin{frame}[shrink=25]{Exemplo: efeito da ordem de inserção}
\begin{columns}[T]
\begin{column}{0.45\textwidth}
\centering\textbf{Aleatória} \{5,2,7,1,3,6,8\}\\[2pt]
\begin{tikzpicture}[level distance=8mm, sibling distance=16mm,
  level 2/.style={sibling distance=8mm}]
  \node[bstnode]{5}
    child{ node[bstnode]{2}
      child{ node[bstnode]{1}} child{ node[bstnode]{3}} }
    child{ node[bstnode]{7}
      child{ node[bstnode]{6}} child{ node[bstnode]{8}} };
\end{tikzpicture}\\[2pt]
\small $h = 2$, balanceada
\end{column}
\begin{column}{0.45\textwidth}
\centering\textbf{Ordenada} \{1,2,3,4,5,6,7\}\\[2pt]
\begin{tikzpicture}[level distance=6mm, sibling distance=8mm]
  \node[bstnode]{1}
    child[missing]
    child{ node[bstnode]{2}
      child[missing]
      child{ node[bstnode]{3}
        child[missing]
        child{ node[bstnode]{4}
          child[missing]
          child{ node[bstnode]{5}
            child[missing]
            child{ node[bstnode]{6}
              child[missing]
              child{ node[bstnode]{7}} } } } } };
\end{tikzpicture}\\[2pt]
\small $h = 6$, degenerada
\end{column}
\end{columns}
\vspace{4pt}
\begin{block}{Motivação para a próxima aula}
\emph{Árvores balanceadas} (AVL, rubro-negra) garantem $h = O(\log n)$ \alert{no pior caso}.
\end{block}
\end{frame}

\begin{frame}[shrink=25]{Demonstração: a árvore vira uma lista}
\centering
Inserindo 1, 2, 3, 4, 5 \emph{em ordem}, a cada passo:\\[6pt]
\begin{tikzpicture}[
  level distance=6mm, sibling distance=8mm,
  every node/.style={font=\scriptsize}
]
  \node[bstnode]{1}; \node at (0,-2.2) {após 1};
\end{tikzpicture}
\hfill
\begin{tikzpicture}[level distance=6mm, sibling distance=8mm,
  every node/.style={font=\scriptsize}]
  \node[bstnode]{1}
    child[missing] child{ node[bstnode]{2}};
  \node at (0.2,-2.2) {após 1,2};
\end{tikzpicture}
\hfill
\begin{tikzpicture}[level distance=6mm, sibling distance=8mm,
  every node/.style={font=\scriptsize}]
  \node[bstnode]{1}
    child[missing]
    child{ node[bstnode]{2}
      child[missing] child{ node[bstnode]{3}}};
  \node at (0.4,-2.2) {após 1,2,3};
\end{tikzpicture}
\hfill
\begin{tikzpicture}[level distance=6mm, sibling distance=8mm,
  every node/.style={font=\scriptsize}]
  \node[bstnode]{1}
    child[missing]
    child{ node[bstnode]{2}
      child[missing]
      child{ node[bstnode]{3}
        child[missing] child{ node[bstnode]{4}}}};
  \node at (0.6,-2.2) {após 1,2,3,4};
\end{tikzpicture}
\hfill
\begin{tikzpicture}[level distance=6mm, sibling distance=8mm,
  every node/.style={font=\scriptsize}]
  \node[bstnode]{1}
    child[missing]
    child{ node[bstnode]{2}
      child[missing]
      child{ node[bstnode]{3}
        child[missing]
        child{ node[bstnode,hl]{4}
          child[missing] child{ node[bstnode,hl]{5}}}}};
  \node at (0.8,-2.2) {após 1,2,3,4,5};
\end{tikzpicture}\\[10pt]
\begin{exampleblock}{Observe}
A árvore degenera em uma \alert{lista encadeada}: $h = n-1$, busca $\Theta(n)$.
O mesmo conjunto de chaves em ordem aleatória dá $h \approx 2$. A \emph{ordem de inserção} é tudo.
\end{exampleblock}
\end{frame}

\begin{frame}[fragile,shrink=10]{Na vida real: \texttt{java.util.TreeMap} / \texttt{TreeSet}}
\begin{block}{O Java já implementa BST balanceada}
\texttt{TreeMap<K,V>} e \texttt{TreeSet<E>} são \emph{árvores rubro-negras}: garantem $O(\log n)$ no pior caso e ordenação automática das chaves.
\end{block}
\begin{lstlisting}
import java.util.TreeMap;
import java.util.NavigableMap;

NavigableMap<String, Integer> agenda = new TreeMap<>();
agenda.put("Ana",    11);
agenda.put("Bruno",  55);
agenda.put("Carla",  99);

agenda.firstKey();          // "Ana"            -- min
agenda.lastKey();           // "Carla"          -- max
agenda.higherKey("Bruno");  // "Carla"          -- sucessor estrito
agenda.headMap("Carla");    // {Ana=11, Bruno=55} -- range query
\end{lstlisting}
\begin{exampleblock}{Quando reimplementar?}
\emph{Quase nunca em produção} --- use \texttt{TreeMap}. Implementamos BST do zero para \alert{entender} a estrutura e poder modificá-la (ex.: AVL, treap, BST aumentada para intervalos).
\end{exampleblock}
\end{frame}

% ============================================================
\section{Exercícios em sala}

\begin{frame}[shrink=20]{Exercício 1 --- Construção e travessias}
\begin{block}{Enunciado}
Insira na BST inicialmente vazia, \emph{nesta ordem}, as chaves:
\[
50,\; 30,\; 70,\; 20,\; 40,\; 60,\; 80,\; 35,\; 45
\]
\textbf{Pedidos:}
\begin{enumerate}
  \item Desenhe a BST resultante.
  \item Liste as chaves em \emph{in-order}, \emph{pre-order} e \emph{post-order}.
  \item Liste a ordem da travessia \emph{BFS} (nível a nível).
  \item Qual é a altura $h$ da árvore?
\end{enumerate}
\end{block}
\begin{exampleblock}{Dica}
A travessia in-order deve produzir as chaves em ordem crescente --- use isso para conferir seu desenho.
\end{exampleblock}
\end{frame}

\begin{frame}[shrink=20]{Exercício 2 --- Busca, sucessor e remoção}
\begin{block}{Use a árvore do Exercício 1}
\begin{enumerate}
  \item Quantas comparações são feitas ao buscar a chave 45? E a chave 55?
  \item Qual é o sucessor in-order de 40? E o sucessor de 50?
  \item Remova a chave 30 (nó com dois filhos) usando o método de Hibbard. Desenhe a árvore resultante.
  \item Remova a chave 70. Desenhe a árvore resultante.
  \item A altura mudou? Como?
\end{enumerate}
\end{block}
\end{frame}

\begin{frame}{Exercício 3 --- Pior caso e implementação}
\begin{block}{Parte A --- pior caso}
Insira em uma BST vazia, \emph{nesta ordem}: $1, 2, 3, 4, 5, 6, 7$.
\begin{enumerate}
  \item Desenhe a árvore.
  \item Quantas comparações são necessárias para buscar a chave 7?
  \item Quanto custaria a mesma busca em uma BST balanceada com as mesmas chaves?
\end{enumerate}
\end{block}
\begin{block}{Parte B --- código}
Escreva, em Java, um método \texttt{int contarFolhas(Node<K,V> x)} que retorna o número de folhas de uma subárvore.\\
\emph{Resposta esperada:} 4 funções de uma linha + caso-base. Discuta a complexidade.
\end{block}
\end{frame}

% ============================================================
\section{Recapitulação}

\begin{frame}[shrink=10]{Recapitulando}
\begin{block}{O que aprendemos hoje}
\begin{itemize}\small
  \item \textbf{Árvore binária:} terminologia (raiz, folha, altura), tipos (cheia, completa, perfeita) e duas representações (ligada e por array).
  \item \textbf{Travessias DFS:} in-order (ordenado), pre-order, post-order; todas em $\Theta(n)$.
  \item \textbf{Travessia BFS:} nível a nível, usando fila (\texttt{ArrayDeque}).
  \item \textbf{Propriedade da BST:} esquerda $\le$ raiz $\le$ direita.
  \item \textbf{Search, Min, Max, Successor, Insert, Delete} custam $O(h)$.
  \item \textbf{Inserção (Sedgewick)} recursiva retornando a subárvore; \textbf{Remoção (Hibbard)} usa sucessor in-order.
  \item \textbf{Altura} é tudo: $O(\log n)$ esperado, $\Theta(n)$ no pior caso $\Rightarrow$ motiva árvores balanceadas.
\end{itemize}
\end{block}
\begin{exampleblock}{Próxima aula}
Árvores balanceadas (AVL): rotações que mantêm $h = O(\log n)$ no pior caso.
\end{exampleblock}
\end{frame}

% ============================================================
\section*{Referências}
\begin{frame}[allowframebreaks]{Referências}
\nocite{*}
\bibliographystyle{plain}
\footnotesize
\bibliography{../referencias}
\end{frame}

\end{document}
